\documentclass[11pt]{article}
\usepackage[a4paper, top=18mm, bottom=18mm, left=20mm, right=20mm]{geometry}
\usepackage[T2A]{fontenc}
\usepackage[utf8]{inputenc}
\usepackage[ukrainian]{babel}
\usepackage{graphicx}
\usepackage[export]{adjustbox}
\usepackage{amsmath}
\usepackage{amssymb}
\graphicspath{{/app/Quiz/Scripts/2020/Mod2020_4/BinTree/}{/app/Quiz/Scripts/2021/Mod2021_3/BinTree/}}
\setlength{\parindent}{0pt}
\setlength{\parskip}{6pt}
\emergencystretch=3em
\pagestyle{empty}
\begin{document}
%begin
{\centering\Large\bfseries Контрольна робота\par}
\medskip
\noindent\rule{\linewidth}{0.5pt}
\smallskip
\noindent\textbf{Варіант \No\,1}\hfill \textit{ПІБ:}\,\underline{\hspace{60mm}}
\vspace{4mm}

%beginitem
1. \textbf{Що буде виведено за виконання фрагменту коду?}
%code
#include <iostream>
using namespace std;

int f(int &x, int &y){
    x = 1;
    y+= 4;
    return x;
}

int main(){
    int a = 9, b = 5;
    a = f(a, a);
    cout << a << ":" << b <<':';
    {
        int a = 7, b = 2;
        cout << ((a<3) && ((b+=1) < 3)) << ':';
        cout << a << ":" << b << ':';
    }
    {
        inta = 8;
        cout << a << ':';
    }
    cout << a << ":" << b << endl;
    return 0;
}
%code

\par\vspace{5mm}
%enditem

%beginitem
2. \textbf{Що буде виведено за виконання фрагменту коду?}
%code
#include <iostream>
using namespace std;

int f(int &x, int &y){
    x = 1;
    y+= 4;
    return x;
}

int main(){
    int a = 9, b = 5;
    a = f(a, a);
    cout << a << ":" << b <<':';
    {
        int a = 7, b = 2;
        cout << ((a<3) && ((b+=1) < 3))
             << a << ":" << b << ":";
    }
    {
        inta = 8;
        cout << a << ':';
    }
    cout << a << ":" << b << endl;
    return 0;
}
%code

\par\vspace{5mm}
%enditem

%beginitem
3. 
Вкажіть значення та тип виразу:
1<"4"
\par\vspace{5mm}
%enditem

%beginitem
4. %goodpath:Scripts/2018/Mod2018_1/01-good.txt
\textbf{Відмітити все, що є однією правильною лексемою:}
( 1 ) &&

( 2 ) 923q

( 3 ) ||

( 4 ) ()
\par\vspace{5mm}
%enditem

%beginitem
5. \textbf{Що буде виведено за виконання фрагменту коду?}

%code
print(2 < 2.0 >= 5.0 > 5)
%code
\par\vspace{5mm}
%enditem

%beginitem
6. 
Кожен рядок текстового файлу містить по два дійсних числа, розділені білими символами.
Написати функцію, що для заданого текстового файлу будує новий текстовий файл, рядки якого крім зображень дйсних чисел ще містять результат обчислення на них функції $f$. У вихідному файлі зображення чисел мають розділятися пробілами. 
$$f(x, y) =\frac{\ln x + 50}{(\sin x + 0.1) (y + 49)}$$ 
Якість коду також оцінюється.

\par\vspace{5mm}
%enditem

%beginitem
7. \textbf{Що буде виведено за виконання фрагменту коду?}
try:
    res = -1**True+-2
    if isinstance(res, bool):
        print(f'bool;{res}')
    elif isinstance(res, int):
        print(f'int;{res}')
    elif isinstance(res, float):
        print(f'float;{res:.2f}')
    else:
        print('???')
except:
    print('ERR')
\par\vspace{5mm}
%enditem

%beginitem
8. \textbf{Що буде виведено за виконання фрагменту коду?}

%code
#include <iostream>
using namespace std;

bool f(int n){
    cout<<"f";
    return n;
}

int main(){
    cout<<(f(1) and f(-6));
    return 0;
}
%code
\par\vspace{5mm}
%enditem

%beginitem
9. \textbf{Що буде виведено за виконання фрагменту коду?}
try:
    i = 15
    while i>0:
        i += -3
    print(i, end=';')
except:
    print('Z')
\par\vspace{5mm}
%enditem

%beginitem
10. \textbf{Що буде виведено за виконання фрагменту коду?}

def f(a):
    u=47
    if a<5: 
        return 0
    if a>=-2:
         u=3
    else:
         u=8
    return u

print(f(1))
\par\vspace{5mm}
%enditem

%beginitem
11. \textbf{Що буде виведено за виконання фрагменту коду?}

x,y,z=0,3,1
x,x,y,y=x,y,6,z
print(x,y,z)
\par\vspace{5mm}
%enditem

%beginitem
12. 

\begin{center}
	\adjustimage{max size={0.8\linewidth}{0.8\paperheight}}
	{../BinTree/trees_exp/37.png}
\end{center}
Указати послідовність міток вершин дерева, зображеного на рис., що утвориться в результаті обходу дерева в прямому порядку. Мітки записувати через пробіл.\par Обхід дерева починається з кореня (на рис. це верхня вершина).
%answer:37 прямому
\par\vspace{5mm}
%enditem

%beginitem
13. \textbf{Що буде виведено за виконання фрагменту коду?}
%code
if (0 < 0)
    cout << "a";
else
    cout << "s";
%code

\par\vspace{5mm}
%enditem

%beginitem
14. \textbf{Що буде виведено за виконання фрагменту коду?}

%code
1<"4"
%code
\par\vspace{5mm}
%enditem

%beginitem
15. \textbf{Що буде виведено за виконання фрагменту коду?}

%code
print(3/3*3*3)
%code
\par\vspace{5mm}
%enditem

%beginitem
16. \textbf{Що буде виведено за виконання фрагменту коду?}
%code

x,y,z=0,3,1
x,x,y,y=x,y,6,z
print(x,y,z)
%code
\par\vspace{5mm}
%enditem

%beginitem
17. 
Написати програму, що запитує у користувача значення дійсних параметрів $x$ та $y$, обчислює для них значення виразу 
$$\frac{\ln x + 50}{(\sin x + 0.1) (y + 49)}$$ 
та повiдомляє введенi данi та результат обчислення.
 Оцінюється правильність та якість коду.

\par\vspace{5mm}
%enditem

%beginitem
18. 
Написати програму, що запитує у користувача значення дійсних параметрів $x$ та $y$, обчислює для них значення виразу та повiдомляє введенi данi та результат обчислення.
$$\frac{\ln x + 50}{(\sin x + 0.1) (y + 49)}$$ 

\par\vspace{5mm}
%enditem

%beginitem
19. %goodpath:Scripts/2019/Mod2019_1/Lex/03-good.txt
\textbf{Відмітити всі коректні оператори Python:}
( 1 ) x<y and y>0

( 2 ) y=3^^2;

( 3 ) x<3 or y<7

( 4 ) print(end="1","qwe")
\par\vspace{5mm}
%enditem

%beginitem
20. \textbf{Що буде виведено за виконання фрагменту коду?}

%code
a = ('a','b','c','d','e',0,1)
print(a[0])
%code

\par\vspace{5mm}
%enditem

%beginitem
21. \textbf{Що буде виведено за виконання фрагменту коду?}

%code
a = 0
b = 3
c = 8

def f(a):
global c    a=1
    b+=4
    c=3
    return a+b+c

a, b, c = 7, 5, 6
try:
    print(f(a), a, b, c, sep=';')
except:
    print('Z', a, b, c, sep=';')
%code
\par\vspace{5mm}
%enditem

%beginitem
22. \textbf{Що буде виведено за виконання фрагменту коду?}
try:
    res = not2<2.0 and not5.0>=5 and True<=3
    if isinstance(res, bool):
        print(f'bool;{res}')
    elif isinstance(res, int):
        print(f'int;{res}')
    elif isinstance(res, float):
        print(f'float;{res:.2f}')
    else:
        print('???')
except:
    print('Z')
\par\vspace{5mm}
%enditem

%beginitem
23. \textbf{Що буде виведено за виконання фрагменту коду?}

%code
print(1 < "4")
%code
\par\vspace{5mm}
%enditem

%beginitem
24. \textbf{Що буде виведено за виконання фрагменту коду?}

print('R', end=' ')
a=[10,11,12,13,14]
a.remove('10')
print(a)

\par\vspace{5mm}
%enditem

%beginitem
25. \textbf{Що буде виведено за виконання фрагменту коду?}

print('R', end=' ')
for a in range(-5,-5,-1):
    if a<3:
        continue
        print(a, end=' ')
    if a<3:
        break
    else:
        print('end', end=' ')
    print(a)
\par\vspace{5mm}
%enditem

%beginitem
26. \textbf{Що буде виведено за виконання фрагменту коду?}

%code
i=15
while i>0:
    i+=-3
print(i, end=' ')
%code

\par\vspace{5mm}
%enditem

%beginitem
27. \textbf{Що буде виведено за виконання фрагменту коду?}

print('R', end=' ')
a='0123456789'
print(a[-5:-5:-1])

\par\vspace{5mm}
%enditem

%beginitem
28. \textbf{Що буде виведено за виконання фрагменту коду?}

%code
a,b,c=0,3,8
def f(a):
global c    a=1
    b+=4
    c=3
    return a+b+c

a,b,c=7,5,6
print(f(a),a,b,c)
%code
\par\vspace{5mm}
%enditem

%beginitem
29. \textbf{Що буде виведено за виконання фрагменту коду?}

%code
-1**True+-2
%code
\par\vspace{5mm}
%enditem

%beginitem
30. \textbf{Що буде виведено за виконання фрагменту коду?}

%code
try:
    i = 15
    while i>0:
        i += -3
    print(i, end=';')
except:
    print('Z')

%code
\par\vspace{5mm}
%enditem

%beginitem
31. \textbf{Що буде виведено за виконання фрагменту коду?}
try:
    i = 0
    while i<5:
        i += 1
    print(i, end=';')
except:
    print('ERR')
\par\vspace{5mm}
%enditem

%beginitem
32. \textbf{Що буде виведено за виконання фрагменту коду?}

%code
try:
    i = 0
    while i<5:
        i += 1
    print(i, end=';')
except:
    print('Z')

%code
\par\vspace{5mm}
%enditem

%beginitem
33. \textbf{Що буде виведено за виконання фрагменту коду?}

%code
print('R', end=' ')
a='0123456789'
print(a[-5:-5:-1])
%code

\par\vspace{5mm}
%enditem

%beginitem
34. \textbf{Що буде виведено за виконання фрагменту коду?}

%code
cout << 15 * 15 * 15 * 15;
%code
\par\vspace{5mm}
%enditem

%beginitem
35. \textbf{Що буде виведено за виконання фрагменту коду?}
%code
if (0 < 0)
    cout << "a";
else
    cout << "a";
%code

\par\vspace{5mm}
%enditem

%beginitem
36. \textbf{Що буде виведено за виконання фрагменту коду?}
def f(arg, n):
    n = 0
    tmp = arg
    tmp[-9:-9] = [4,5]
    return len(tmp)>3

try:
    a = [10,11,12,13,14]
    n = 3
    f(a, n)
    print(n, end=';')
    for el in a:
        print(el, end=';')
except:
    print('Z')
\par\vspace{5mm}
%enditem

%beginitem
37. \textbf{Що буде виведено за виконання фрагменту коду?}
%code
a,b,c=6,9,7
def f(a,b,c):
    print(a,b,c,end=" ")

f(0,3)
print(a,b,c)
%code
\par\vspace{5mm}
%enditem

%beginitem
38. 

\begin{center}
	\adjustimage{max size={0.8\linewidth}{0.8\paperheight}}
	{../BinTree/trees_exp/37.png}
\end{center}
Указати послідовність міток вершин дерева, зображеного на рис., що утвориться в результаті обходу дерева в прямому порядку. Мітки записувати через пробіл.\par Алгоритм починає роботу з кореня (на рис. це верхня вершина).
%answer:37 прямому
\par\vspace{5mm}
%enditem

%beginitem
39. \textbf{Що буде виведено за виконання фрагменту коду?}

%code
cout << (!2<2.0 and !5.0>= 5 and true<=3);
%code
\par\vspace{5mm}
%enditem

%beginitem
40. \textbf{Що буде виведено за виконання фрагменту коду?}
%Оригинал был утерян. Требует отладки!
%code
_a = 0

class A:
    _a = 1
    
    def __init__(self):
        self._a = 2
        _a = 3
        A._a = 4

obj = A()
obj._a = 5
try:
    print(_a, end=' ')
    print(obj._a, end=' ')
    print(A._a, end=' ')
    print(obj._a, end=' ')
    print(obj._A__a)
except:
    print('ERR')
%code
\par\vspace{5mm}
%enditem

%beginitem
41. \textbf{Що буде виведено за виконання фрагменту коду?}
code
a = 0
h = 1

class A:
    a = 2
    
    def __init__(self):
global a        self.a = 3
        a = 4
        A.a = 5

obj = A()
print(a, h, A.a, obj.a, sep=' ')
code
\par\vspace{5mm}
%enditem

%beginitem
42. \textbf{Що буде виведено за виконання фрагменту коду?}

a,b,c=0,3,8
def f(a):
global c    a=1
    b=4
    c=3
    return a+b+c

a,b,c=7,5,6
print(f(a),a,b,c)
\par\vspace{5mm}
%enditem

%beginitem
43. 
Записати ВИРАЗ, значенням якого є множина, що складається з цілих чисел від 20 до 2010 (включно), що не діляться на 2.\par\vspace{5mm}
%enditem

%beginitem
44. \textbf{Що буде виведено за виконання фрагменту коду?}

%code
#include <iostream>

int f(int a, int b){
    int c = 47;
    if (a < 5)
        return 0;
    if (b >= -2)
         c = 3;
    else 
        c = 8;
    return c;
}

int main(){
    std::cout << f(1, 1);
    return 0;
}
%code
\par\vspace{5mm}
%enditem

%beginitem
45. \textbf{Що буде виведено за виконання фрагменту коду?}

%code
print(-1**True+-2)
%code
\par\vspace{5mm}
%enditem

%beginitem
46. 
Змінні \kw{next} та \kw{prev} вузла списку позначають наступний та попередній за ним вузол відповідно. Починаючи з голови, у список записано послідовні цілі числа від~$-33$ до~$15$. Нехай змінна \kw{p} позначає вузол, що зберігає значення~$1$.
Яке значення зберігається у вузлі \kw{p.prev.prev.prev.prev.prev} ?
\par\vspace{5mm}
%enditem

%beginitem
47. \textbf{Що буде виведено за виконання фрагменту коду?}
%code
for a in range(-5,-5,-1):
    if a<3:
        continue
        print(a,end=' ')
    if a<3:
        break
else:
    print('end', end=' ')
print(a, end=' ')
%code
\par\vspace{5mm}
%enditem

%beginitem
48. \textbf{Що буде виведено за виконання фрагменту коду?}

print('R', end=' ')
i=15
while i>0:
    i+=-3
print(i)

\par\vspace{5mm}
%enditem

%beginitem
49. \textbf{Що буде виведено за виконання фрагменту коду?}

%code
print('R', end=' ')
a=[10,11,12,13,14]
a.remove('10')
print(a)
%code

\par\vspace{5mm}
%enditem

%beginitem
50. \textbf{Що буде виведено за виконання фрагменту коду?}
try:
    res = 3/3*3*3
    if isinstance(res, bool):
        print(f'bool;{res}')
    elif isinstance(res, int):
        print(f'int;{res}')
    elif isinstance(res, float):
        print(f'float;{res:.2f}')
    else:
        print('???')
except:
    print('ERR')
\par\vspace{5mm}
%enditem

%beginitem
51. \textbf{Що буде виведено за виконання фрагменту коду?}
try:
    for a in range(-5, -5, -1):
        if a<3:
            continue
        print(a, end=';')
        if a<3:
            break
    else:
        print('end',end=';')
    print(a)
except:
    print('Z')
\par\vspace{5mm}
%enditem

%beginitem
52. \textbf{Що буде виведено за виконання фрагменту коду?}

try:
    print(0,end="")
    print(3/1,end="")
    print(8,end="")
except Exception: print(7,end="")
except TypeError: print(5,end="")
else: print(6,end="")
finally: print(2,end="")
\par\vspace{5mm}
%enditem

%beginitem
53. \textbf{Що буде виведено за виконання фрагменту коду?}
try:
    res = 1<"4j"
    if isinstance(res, bool):
        print(f'bool;{res}')
    elif isinstance(res, int):
        print(f'int;{res}')
    elif isinstance(res, float):
        print(f'float;{res:.2f}')
    else:
        print('???')
except:
    print('Z')
\par\vspace{5mm}
%enditem

%beginitem
54. %goodpath:Scripts/2018/Mod2018_1/03-good.txt
\textbf{Відмітити все, що є коректним оператором С++:}
( 1 ) (x<3 || y<7)

( 2 ) int n; bool n;

( 3 ) x<y<z

( 4 ) int x*y=3;
\par\vspace{5mm}
%enditem

%beginitem
55. \textbf{Що буде виведено за виконання фрагменту коду?}

%code
cout << (15 * 15 * 15 * 15);
%code
\par\vspace{5mm}
%enditem

%beginitem
56. 
Змінні $next$ та $prev$ вузла списку позначають наступний та попередній за ним вузол відповідно. Починаючи з голови, у список записано послідовні цілі числа від $-33$ до $15$. Нехай змінна $p$ позначає вузол, що зберігає значення $1$.
Яке значення зберігається у вузлі $p.prev.prev.prev.prev.prev$ ?
\par\vspace{5mm}
%enditem

%beginitem
57. \textbf{Що буде виведено за виконання фрагменту коду?}

%code
print('R', end=' ')
i=0
while i<5:
    i+=1
print(i)
%code

\par\vspace{5mm}
%enditem

%beginitem
58. 
Задано тип вузла списку
struct Node \{int n; Node *prev, *next;\};
У списку зберігаються послідовні цілі числа від -33 до 15. Вказівник p встановлено на вузол, в якому зберігається число 1.
Чому дорівнює значення p->prev->prev->prev->prev->prev->n ?\par\vspace{5mm}
%enditem

%beginitem
59. \textbf{Що буде виведено за виконання фрагменту коду?}

%code
print(not2<2.0 and not5.0>=5 and True<=3)
%code
\par\vspace{5mm}
%enditem

%beginitem
60. \textbf{Що буде виведено за виконання фрагменту коду?}

%code
#include <iostream>
using namespace std;

int f(int a){
    int u = 47;
    if (a < 5) 
        return 0;
    if (a >= -2)
         u = 3;
    else
         u = 8;
    return u;
}

int main(){
    cout << f(1);
    return 0;
}
%code

\par\vspace{5mm}
%enditem

%beginitem
61. \textbf{Що буде виведено за виконання фрагменту коду?}
%code

def f(n):
    print("f", end="");
    return n

print(f(1) and f(-6))
%code
\par\vspace{5mm}
%enditem

%beginitem
62. \textbf{Що буде виведено за виконання фрагменту коду?}

%code
print('R', end=' ')
for a in range(0,0,-1):
    if a<3:
        continue
        print(a, end=' ')
    if a<3:
        break
    else:
        print('end', end=' ')
    print(a, end=' ')
%code
\par\vspace{5mm}
%enditem

%beginitem
63. \textbf{Що буде виведено за виконання фрагменту коду?}
try:
    for a in range(-5, -5, -1):
        if a<3:
            continue
        print(a, end=';')
        if a<3:
            break
    else:
        print('end',end=';')
    print(a)
except:
    print('ERR')
\par\vspace{5mm}
%enditem

%beginitem
64. \textbf{Що буде виведено за виконання фрагменту коду?}
a = 0
b = 3
c = 8

def f(a):
global c    a = 1
    b = 4
    c = 3
    return a + b + c

a = 7
b = 5
c = 6

try:
    print(f(a), a, b, c, sep=';')
except:
    print('Z', a, b, c, sep=';')
\par\vspace{5mm}
%enditem

%beginitem
65. \textbf{Що буде виведено за виконання фрагменту коду?}

%code
a=[10,11,12,13]
a.remove('10')
print(a)
%code

\par\vspace{5mm}
%enditem

%beginitem
66. \textbf{Що буде виведено за виконання фрагменту коду?}

%code
i=0
while i<5:
    i+=1
print(i, end=' ')
%code

\par\vspace{5mm}
%enditem

%beginitem
67. 
$a$ є цілочисловим масивом типу $numpy.ndarray$ розмірів $2{\times}2$. Сума елементів масиву $a$ дорівнює 30. Чому дорівнює сума елементів масиву $a$ після виконання 
\begin{verbatim}
a += 30
\end{verbatim}


Якщо код некоректний, то в якості відповіді зазначити ERROR.


\par\vspace{5mm}
%enditem

%beginitem
68. \textbf{Що буде виведено за виконання фрагменту коду?}
try:
    res = -1**True+-2
    if isinstance(res, bool):
        print(f'bool;{res}')
    elif isinstance(res, int):
        print(f'int;{res}')
    elif isinstance(res, float):
        print(f'float;{res:.2f}')
    else:
        print('???')
except:
    print('Z')
\par\vspace{5mm}
%enditem

%beginitem
69. 
В умовах попередньої задачі було виконано p->prev=p->next->prev;

Чому дорівнює значення p->prev->prev->prev->prev->n ?\par\vspace{5mm}
%enditem

%beginitem
70. %goodpath:Scripts/2019/Mod2019_1/Lex/01-good.txt
\textbf{Відмітити все, що є однією правильною лексемою:}
( 1 ) %

( 2 ) 1_2_

( 3 ) **

( 4 ) -56
\par\vspace{5mm}
%enditem

%beginitem
71. \textbf{Що буде виведено за виконання фрагменту коду?}

%code
#include <iostream>

bool f(int n){
    std::cout<<"f";
    return n;
}

int main(){
    std::cout<<(f(1) and f(-6));
    return 0;
}
%code
\par\vspace{5mm}
%enditem

%beginitem
72. 
Написати програму, що запитує у користувача значення дійсних параметрів $x$ та $y$, обчислює для них значення виразу 
$$\frac{\ln x + 50}{(\sin x + 0.1) (y + 49)}$$ 
та повiдомляє введенi данi та результат обчислення.

\par\vspace{5mm}
%enditem

%beginitem
73. \textbf{Що буде виведено за виконання фрагменту коду?}

%code
#include <iostream>
using namespace std;

int f(int a, int b){
    int c = 47;
    if (a < 5)
        return 0;
    if (b >= -2)
         c = 3;
    else 
        c = 8;
    return c;
}

int main(){
    cout << f(1, 1);
    return 0;
}
%code
\par\vspace{5mm}
%enditem

%beginitem
74. \textbf{Що буде виведено за виконання фрагменту коду?}
try:
    res = 3/3*3*3
    if isinstance(res, bool):
        print(f'bool;{res}')
    elif isinstance(res, int):
        print(f'int;{res}')
    elif isinstance(res, float):
        print(f'float;{res:.2f}')
    else:
        print('???')
except:
    print('Z')
\par\vspace{5mm}
%enditem

%beginitem
75. \textbf{Що буде виведено за виконання фрагменту коду?}

%code
for a in range(0,0+2,-3):
    if a<3:
        continue
    print(a, end=' ')
    a=0
    if a<3:
        break
else:
    print(13, end=' ')
print(a, end=' ')
%code
\par\vspace{5mm}
%enditem

%beginitem
76. \textbf{Що буде виведено за виконання фрагменту коду?}

def f(a,b):
    c=47
    if a<5:
        return 0
    if b>=-2:
         c=3
    else: 
        c=8
    return c

print(f(1,1))\par\vspace{5mm}
%enditem

%beginitem
77. \textbf{Що буде виведено за виконання фрагменту коду?}

print('R', end=' ')
i=0
while i<5:
    i+=1
print(i)

\par\vspace{5mm}
%enditem

%beginitem
78. \textbf{Що буде виведено за виконання фрагменту коду?}

%code
def f(a,b):
    c=47
    if a<5:
        return 0
    if b>=-2:
         c=3
    else: 
        c=8
    return c

print(f(1,1))
%code
\par\vspace{5mm}
%enditem

%beginitem
79. \textbf{Що буде виведено за виконання фрагменту коду?}

%code
cout << (3 * 3 * 3 * 3);
%code
\par\vspace{5mm}
%enditem

%beginitem
80. \textbf{Що буде виведено за виконання фрагменту коду?}

a,b,c=6,9,7
def f(a,b,c):
    print(a,b,c,end="")

f(0,3)
print(a,b,c)

\par\vspace{5mm}
%enditem

%beginitem
81. 
Вкажіть значення та тип виразу:
3/3*3*3
\par\vspace{5mm}
%enditem

%beginitem
82. \textbf{Що буде виведено за виконання фрагменту коду?}

%code
a,b,c=0,3,8
def f(a):
    a=1
    b+=4
    c=3
    return a+b+c

a,b,c=7,5,6
print(f(a),a,b,c)
%code
\par\vspace{5mm}
%enditem

%beginitem
83. \textbf{Що буде виведено за виконання фрагменту коду?}
try:
    res = 1<"4j"
    if isinstance(res, bool):
        print(f'bool;{res}')
    elif isinstance(res, int):
        print(f'int;{res}')
    elif isinstance(res, float):
        print(f'float;{res:.2f}')
    else:
        print('???')
except:
    print('ERR')
\par\vspace{5mm}
%enditem

%beginitem
84. \textbf{Що буде виведено за виконання фрагменту коду?}
a=[10,11,12,13,14]; a[-9:-9]=(); print(a)\par\vspace{5mm}
%enditem

%beginitem
85. 
Значенням змінної \kw{a} є цілочисловий масив типу $numpy.ndarray$ розмірів $3{\times}3$. На скільки збільшиться сума елементів масиву \kw{a} після виконання
\begin{verbatim}
a[-3, -4] += 1
\end{verbatim}


Якщо код некоректний, то в якості відповіді зазначити ERROR.


\par\vspace{5mm}
%enditem

%beginitem
86. \textbf{Що буде виведено за виконання фрагменту коду?}
a = 0
b = 3
c = 8

def f(a):
global c    a = 1
    b += 4
    c = 3
    return a + b + c

a = 7
b = 5
c = 6

try:
    print(f(a), a, b, c, sep=';')
except:
    print('Z', a, b, c, sep=';')
\par\vspace{5mm}
%enditem

%beginitem
87. 
Змінні \kw{next} та \kw{prev} вузла списку позначають наступний та попередній за ним вузол відповідно. Починаючи з голови, у список записано послідовні цілі числа від~$-33$ до~$15$. Нехай змінна \kw{p} позначає вузол, що зберігає значення~$1$.
Яке число зберігається у вузлі \kw{p.prev.prev.prev.prev.prev} ?
\par\vspace{5mm}
%enditem

%beginitem
88. \textbf{Що буде виведено за виконання фрагменту коду?}
def f(arg, n):
    n = 0
    tmp = arg
    tmp[-9:-9] = [4,5]
    return len(tmp)>3

try:
    a = [10,11,12,13,14]
    n = 3
    f(a, n)
    print(n, end=';')
    for el in a:
        print(el, end=';')
except:
    print('ERR')
\par\vspace{5mm}
%enditem

%beginitem
89. \textbf{Що буде виведено за виконання фрагменту коду?}

%code
a=[10,11,12,13]
a[-9:-9]=[4,5]
print(a)
%code

\par\vspace{5mm}
%enditem

%beginitem
90. \textbf{Що буде виведено за виконання фрагменту коду?}

%code
a,b,c=0,3,8
def f(a):
    a=1
    b=4
    c=3
    return a+b+c

a,b,c=7,5,6
print(f(a),a,b,c)
%code
\par\vspace{5mm}
%enditem

%beginitem
91. \textbf{Що буде виведено за виконання фрагменту коду?}

%code
print('R', end=' ')
a=[10,11,12,13,14]
a[-9:-9]=()
print(a)
%code

\par\vspace{5mm}
%enditem

%beginitem
92. \textbf{Що буде виведено за виконання фрагменту коду?}
def f(n):
    print('f', end='')
    return n

try:
    print(f(1) and f(-6))
except:
    print('ERR')
\par\vspace{5mm}
%enditem

%beginitem
93. \textbf{Що буде виведено за виконання фрагменту коду?}

%code
cout << 3 * 3 * 3 * 3;
%code
\par\vspace{5mm}
%enditem

%beginitem
94. \textbf{Що буде виведено за виконання фрагменту коду?}

%code
#include <iostream>

int f(int a){
    int u = 47;
    if (a < 5) 
        return 0;
    if (a >= -2)
         u = 3;
    else
         u = 8;
    return u;
}

int main(){
    std::cout << f(1);
    return 0;
}
%code

\par\vspace{5mm}
%enditem

%beginitem
95. \textbf{Що буде виведено за виконання фрагменту коду?}
try:
    res = not2<2.0 and not5.0>=5 and True<=3
    if isinstance(res, bool):
        print(f'bool;{res}')
    elif isinstance(res, int):
        print(f'int;{res}')
    elif isinstance(res, float):
        print(f'float;{res:.2f}')
    else:
        print('???')
except:
    print('ERR')
\par\vspace{5mm}
%enditem

%beginitem
96. 
Записати ВИРАЗ, значенням якого є список, що складається з цілих чисел від 20 до 2010 (включно), причому спочатку за зростанням записано всі, що не діляться на 2, а потім решту за зростанням.\par\vspace{5mm}
%enditem

%beginitem
97. \textbf{Що буде виведено за виконання фрагменту коду?}

%code
#include <iostream>
using namespace std;

int a = 0, b = 3, c = 8;

int f(int &a){
int c;    a += 1;
    b = 4;
    c = 3;
    return a + b + c;
}

int main(){
    inta = 7;
    intb = 5;
    intc = 6;
    cout << f(a) << ':';
    cout << a << ':' << b << ':' << c;
    return 0; 
}
%code
\par\vspace{5mm}
%enditem

%beginitem
98. \textbf{Що буде виведено за виконання фрагменту коду?}
try:
    print(0, end=';')
    print(3/1, end=';')
    print(8, end=';')
except Exception:
    print(7, end=';')
except ZeroDivisionError:
    print(5, end=';')
except:
    print('Z', end=';')
else:
    print(6, end=';')
finally:
    print(2)
\par\vspace{5mm}
%enditem

%beginitem
99. \textbf{Що буде виведено за виконання фрагменту коду?}
for a in range(-5,-5,-1):
    if a<3:
        continue
        print(a,end=' ')
    if a<3:
        break
else:
    print('end', end=' ')
print(a, end=' ')
\par\vspace{5mm}
%enditem

%beginitem
100. \textbf{Що буде виведено за виконання фрагменту коду?}

%code
not2<2.0 and not5.0>=5 and True<=3
%code
\par\vspace{5mm}
%enditem

%end
\newpage
\end{document}

